\naslov{Centralni limitni izrek}

\begin{izrek}[centralni limitni izrek]
  Naj bodo $X_1, X_2, \ldots$ neodvisne in enako porazdeljene slučajne
  spremenljivke s končnim drugim momentom.
  Označimo $\mu_1 = E(X_1)$, $\sigma_1^2 = \var(X_1)$ in $S_n = X_1 + \cdots +
  X_n$.
  Za $W_n = \frac{S_n - \mu_1 n}{\sigma_1 \sqrt{n}}$ potem velja
  \[
	\lim_{n \to \infty} P(W_n \le w) = \phi(w)
  \]
  enakomerno za $w \in \R$, torej
  \[
	\lim_{n \to \infty} \sup_{w \in \R} \abs{P(W_n \le w) - \phi(w)}
	= 0.
  \]
\end{izrek}

Izrek lahko uporabimo za ocenjevanje porazdelitve vsote veliko IID slučajnih
spremenljivk.
Izračunamo $\mu_1 = E(X_1)$ in $\sigma_1^2 = \var(X_1)$ in ocenimo verjetnost
kot
\[
  P(S_n \le x) = \phi\!\left( \frac{x - n \mu_1}{\sigma_1 \sqrt{n}} \right).
\]
Za $x$ vzamemo mejo, ki nas zanima, pri čemer za vsote z vrednostmi na mreži $a
\Z + b$ za $a, b \in \N$ po dogovoru vzamemo srednjo vrednost; torej za
ocenjevanje verjetnosti, da je padlo manj kot $M$ pik, npr.~vzamemo $x = M -
\pol$, za več kot $M$ pa $x = M + \pol$.

Povemo lahko tudi nekaj o napaki te ocene.
Če so $X_1, X_2, \ldots$ neodvisne, za $\mu_n = E(S_n)$ in $\sigma_n^2 =
\var(S_n)$ velja
\[
  \sup_{x \in \R} \abs{ P(S_n \le x) - \phi\!\left( \frac{x - \mu_n}{\sigma_n}
	\right) }
  \le \frac{0.5583}{\sigma_n^3} \sum_{k=1}^n E\!\left( (X_k - E(X_k))^3 \right).
\]
Pri tem nismo predpostavili, da so spremenljivke enako porazdeljene, torej lahko
to oceno uporabimo v več primerih.
Če npr.~pokažemo, da desna stran konvergira k $0$ za $n \to \infty$, lahko
pokažemo rezultat CLI tudi za različno porazdeljene slučajne spremenljivke.

\naslov{Konvergenca slučajnih spremenljivk}

Zaporedje $X_1, X_2, \ldots$ konvergira proti $X$
\begin{itemize}
\item \pojem{šibko}, če je za vsako zvezno in omejeno $h$
  \[
	\lim_{n \to \infty} E(h(X_n)) = E(h(X)),
  \]
\item \pojem{v verjetnosti}, če je za vsak $\varepsilon > 0$
  \[
	\lim_{n \to \infty} P(d(X_n, X) > \varepsilon) = 0,
  \]
\item \pojem{skoraj gotovo}, če je
  \[
	P\!\left( \{ \lim_{n \to \infty} X_n = X \} \right) = 1.
  \]
\end{itemize}
Tretja točka implicira drugo, druga pa prvo.
Če imajo spremenljivke vrednosti v $\R$, je prva točka ekvivalentna pogoju, da
\[
  \lim_{n \to \infty} P(X_n \le x) = P(X \le x)
\]
za vsak $x$ z $P(X=x) = 0$.

\naslov{Cenilke}

Recimo, da imamo model nekega dogajanja in želimo preveriti, če drži vodo.
Karakteristika modela $y$ je neka lastnost (parameter, \ldots), ki je za ta
model značilna, mi pa je ne poznamo.
Iz opažanj lahko ocenimo vrednost te karakteristika $\hat{y}$, oceni pravimo
\pojem{cenilka}.
Za cenilko definiramo \pojem{pričakovano} ali \pojem{srednjo kvadratično napako}
\[
  \operatorname{MSE}(\hat{y} \such y) = E({(y - \hat{y})}^2)
\]
($\hat{y}$ je slučajna spremenljivka, pisana z malo).
Poleg tega definiramo \pojem{pristranskost}
\[
  \operatorname{Bias}(\hat{y} \such y) = E(\hat{y}) - y.
\]
Cenilka je \pojem{nepristranska}, če je $\operatorname{Bias}(\hat{y} \such y) =
0$.
Pri takih cenilkah je MSE enaka varianci, in lahko definiramo
\[
  \operatorname{SE}(\hat{y}) = \operatorname{RMSE}(\hat{y} \such y) =
  \sqrt{\var(\hat{y})},
\]
v splošnem pa dobimo
\[
  \var(\hat{y}) = \operatorname{MSE}(\hat{y} \such y) - \left(
	\operatorname{Bias}(\hat{y} \such y) \right)^2.
\]

Če podatke dobivamo postopoma, dobimo zaporedje cenilk $(\hat{y}_i)_i$.
Za tako zaporedje pravimo, da je \pojem{šibko dosledno}, če
\[
  \hat{y}_n \xrightarrow[n \to \infty]{d} y.
\]
Zadosten pogoj je, da srednja kvadratna napaka konvergira k $0$, čemur pravimo
\pojem{doslednost}.

Pogosta tema je iskanje \pojem{najboljše nepristranske linearne cenilke} (NNLC).
To je preprosto cenilka, ki ima izmed vseh linearnih cenilk najmanjšo srednjo
kvadratno napako.
Če so $X_1, \ldots, X_n$ nekorelirane z enakimi pričakovanimi vrednostmi in
variancami, je NNLC kar povprečje.

\naslov{Slučajni vektorji}

\newcommand{\sv}[1]{\underline{#1}}
\newcommand{\cov}{\operatorname{cov}}

Za slučajni vektor $\sv{X}$ definiramo \pojem{pričakovano vrednost}
\[
  E(\sv{X}) =
  \begin{bmatrix}
	E(X_1) \\ \vdots \\ E(X_n)
  \end{bmatrix},
\]
za par $\sv{X}, \sv{Y}$ pa \pojem{kovariančno matriko}
\[
  \operatorname{Cov}(\sv{X}, \sv{Y}) =
  \begin{bmatrix}
	\cov(X_1, Y_1) & \cdots & \cov(X_1, Y_n) \\
	\vdots & \ddots & \vdots \\
	\cov(X_n, Y_1) & \cdots & \cov(X_n, Y_n)
  \end{bmatrix}
  = E(\sv{X} \sv{Y}^T) - E(\sv{X}) E(\sv{Y})^T.
\]
Za deterministično matriko $A$ je $E(A \sv{X}) = A E(\sv{X})$, podobno za
slučajno matriko $M$ velja $E(A M) = A E(M)$.
Analogni enakosti dobimo za množenje z desne.
Za deterministični matriki $A, B$ in slučajna vektorja $\sv{X}, \sv{Y}$ je
\[
  \operatorname{Cov}(A \sv{X}, B \sv{Y}) = A \operatorname{Cov}(\sv{X}, \sv{Y})
  B^T.
\]

Dva omembe vredna trika pri obračanju slučajnih matrik sta s sledjo.
Ker je sled linearna preslikava, se lepo obnaša s pričakovano vrednostjo, in je
$E(\sled(M)) = \sled(E(M))$.
Spomnimo se, da je skalar enak svoji sledi, da je sled linearna, $\sled(A+B) =
\sled A + \sled B$, in da lahko ciklično zamenjujemo argumente:
\[
  \sled(ABC) = \sled(BCA) = \sled(CAB).
\]

\naslov{Pridobivanje cenilk}

Pri enostavnem slučajnem vzorčenju imamo populacijo velikosti $N$ in vzorec
velikosti $n$.
Iz populacije izberemo enote $K_1, \ldots, K_n$, $K_i \ne K_j$, tako, da so vse
$n$-terice enako verjetne.
Označimo $X_i = X_{K_i}$.
Imamo nepristranske cenilke
\[
  \hat{\mu} = \inv{n} (X_1 + \ldots + X_n),
\]
za $S = \sum_i (X_i - \hat{\mu})^2$
\begin{gather*}
  \hat{\sigma^2} = \frac{N-1}{N} \frac{S}{n-1}, \\
  \hat{\mu^2} = \hat{\mu}^2 \frac{N-n}{N} \frac{S}{n (n-1)},
\end{gather*}
in napaka
\[
  \hat{\text{SE}^2} = \frac{N-n}{N-1} \frac{\hat{\sigma^2}}{n}.
\]

Pri stratificiranem vzorčenju je populacija razdeljena na $k$ stratumov z
velikostmi $N_1, \ldots, N_k$.
Označimo $w_i = N_i / N$.
Na vsakem stratumu vzorčimo enostavno slučajno.

Poznamo dve pogosti metodi za pridobivanje cenilk.
Prva je \pojem{metoda momentov}, kjer cenilke izpeljemo iz predpisov za momente
(potrebujemo toliko momentov, kolikor je argumentov v porazdelitvi).
Potem poiščemo inverzno preslikavo, s katero izrazimo cenilke z momenti;
\[
  (\hat{y}_1, \hat{y}_2, \ldots, \hat{y}_m) = F^{-1}(\overline{X}, \overline{X^2}, \ldots,
  \overline{X^m}).
\]

Druga metoda je \pojem{metoda največjega verjetja}.
Pri tej metodi iz opažanja $X = (X_1, \ldots, X_n)$ in gostote $f_X(x \such
\theta)$ pridobimo cenilko za $\theta$ kot
\[
  \hat{\theta} = \operatorname{arg\,max} L(\theta \such x)
\]
kjer je $L(\theta \such x) = f_X(x \such \theta)$ verjetje.
V diskretnem primeru zamenjamo $f_X$ z verjetnostjo $P(X = x \such \theta)$.
Pogosto se splača namesto verjetja maksimizirati njegov logaritem, torej
rešiti $\partial_\theta \log L = 0$.

% LocalWords:  stratificiranem stratumov stratumu
